\subsection{Training Data}%
\label{sec:data}

\begin{table*}[htbp]
\caption{\textbf{Camera Pose Estimation} across static and dynamic benchmarks.
The metric AUC is higher-is-better. 
Feed-forward models (DA3, PI3, VGGT) are robust across datasets, but still lag behind the dynamic optimization-based method MegaSaM at strict thresholds on Sintel (\eg, AUC@$3^\circ$ of $16.2$ \vs $22.5$). 
In contrast, dynamic optimization methods (MegaSaM, MonST3R) degrade on wide-baseline scenes (\eg, AUC@$30^\circ$ of $38.1$ \vs $86.4$ on ETH3D). 
Our method, however, substantially advances the state of the art across all scenarios, \eg, improving AUC@$3^\circ$ from $22.5$ to $40.0$ on Sintel, with a $77\%$ relative improvement.}%
\label{tab:camera}
% \vspace{-1em}
\centering
\scriptsize
\setlength{\tabcolsep}{6pt}
\resizebox{\textwidth}{!}{%
\begin{tabular}{c *{6}{cc}}
\toprule
\multirow{3}{*}{\textbf{Method}}
& \multicolumn{6}{c}{\textbf{Static Scenes}}
& \multicolumn{6}{c}{\textbf{Dynamic Scenes}} \\
\cmidrule(lr){2-7}\cmidrule(lr){8-13}
& \multicolumn{2}{c}{\textbf{7 Scenes}}
& \multicolumn{2}{c}{\textbf{NRGBD}}
& \multicolumn{2}{c}{\textbf{ETH3D}}
& \multicolumn{2}{c}{\textbf{DyCheck}}
& \multicolumn{2}{c}{\textbf{Sintel}}
& \multicolumn{2}{c}{\textbf{TUM\text{-}Dynamic}} \\
& AUC@3° & AUC@30°
& AUC@3° & AUC@30°
& AUC@3° & AUC@30°
& AUC@3° & AUC@30°
& AUC@3° & AUC@30°
& AUC@3° & AUC@30° \\
\midrule
{MonST3R} 
& 9.0 & 68.3 & 13.9 & 79.7 & 1.7 & 14.3 
& 11.5 & 45.4 & 4.3 & 45.8 & 7.7 & 48.5 \\
MapAnything & 5.8 & 61.4 & 35.2 & 88.9 & 13.2 & 51.0 & 6.1 & 60.3 & 2.9 & 31.6 & 4.3 & 40.2 \\
{MegaSaM} 
& 10.6 & 71.8 & 17.2 & 83.1 & 5.9 & 38.1 
& 26.8 & 53.1 & 22.5 & 58.3 & 15.4 & 59.0 \\
{VGGT}    
& 10.9 & 74.4 & 81.7 & 97.7 & 18.8 & 62.1 
& 21.0 & 78.7 & 15.0 & 50.0 & 16.6 & 61.2 \\
{PI3}     
& 13.3 & 77.0 & 83.8 & 98.2 & 35.3 & 79.6
& 23.3 & 81.0 & 14.8 & 53.5 & 16.1 & 59.2 \\
{DA3} & 18.7 & 78.2 & 86.4 & 98.4 & 46.1 & 87.0 & 32.1 & 83.9 & 16.2 & 52.7 & 20.8 & 62.7 \\
\midrule
{Ours-1B} 
& \underline{29.6} & \underline{83.1} 
& \underline{89.7} & \underline{98.8} 
& \underline{49.8} & \underline{88.5} 
& \underline{38.4} & \underline{87.3} 
& \underline{35.3} & \underline{73.0} 
& \underline{30.2} & \underline{82.3} \\
{Ours-10B} 
& \textbf{36.4} & \textbf{88.2} 
& \textbf{92.5} & \textbf{99.1} 
& \textbf{56.3} & \textbf{90.4} 
& \textbf{43.7} & \textbf{90.9} 
& \textbf{40.0} & \textbf{79.1} 
& \textbf{36.4} & \textbf{87.5} \\
\bottomrule
\end{tabular}%
}
\vspace{-1em}
\end{table*}








\begin{table*}[htbp]
\caption{\textbf{Depth Estimation} across static and dynamic benchmarks.
The metric \(\delta_{1.25}\) denotes the percentage of predicted depths within a factor of the ground truth (higher is better), and AbsRel is the mean absolute relative error (lower is better).}%
\label{tab:depth}
% \vspace{-1em}
\centering
\tiny
\setlength{\tabcolsep}{6pt}
\resizebox{\textwidth}{!}{%
\begin{tabular}{c *{6}{cc}}
\toprule
\multirow{3}{*}{\textbf{Method}}
& \multicolumn{6}{c}{\textbf{Static Scenes}}
& \multicolumn{6}{c}{\textbf{Dynamic Scenes}} \\
\cmidrule(lr){2-7}\cmidrule(lr){8-13}
& \multicolumn{2}{c}{\textbf{7 Scenes}}
& \multicolumn{2}{c}{\textbf{NRGBD}}
& \multicolumn{2}{c}{\textbf{ETH3D}}
& \multicolumn{2}{c}{\textbf{DyCheck}}
& \multicolumn{2}{c}{\textbf{Sintel}}
& \multicolumn{2}{c}{\textbf{TUM\text{-}Dynamic}} \\
& \(\delta_{1.25}\) & {AbsRel}
& \(\delta_{1.25}\) & {AbsRel}
& \(\delta_{1.25}\) & {AbsRel}
& \(\delta_{1.25}\) & {AbsRel}
& \(\delta_{1.25}\) & {AbsRel}
& \(\delta_{1.25}\) & {AbsRel} \\
\midrule
{MonST3R} & 92.4 & 0.075 & 98.4 & 0.030 & 95.8 & 0.056 & 93.3 & 0.068 & 71.9 & 0.263 & 85.0 & 0.148 \\
MapAnything & 92.9 & 0.070 & 98.7 & 0.022 & 96.3 & 0.035 & 97.0 & 0.049 & 72.5 & 0.251 & 93.1 & 0.052 \\
{MegaSaM} & 93.8 & 0.065 & 96.2 & 0.057 & 94.8 & 0.083 & 97.4 & 0.042 & 74.1 & 0.207 & 92.9 & 0.083 \\
{VGGT}    & 91.9 & 0.073 & 99.1 & 0.019 & 97.4 & 0.036 & 95.2 & 0.055 & 79.2 & 0.189 & 92.2 & 0.064 \\
{PI3}     & 92.8 & 0.068 & 99.2 & 0.011 & {99.6} & 0.016 & 97.4 & 0.041 & 82.5 & 0.144 & 95.5 & 0.046 \\
{DA3} & 93.0 & 0.063 & 99.5 & 0.010 & 99.6 & 0.015 & 97.7 & 0.039 & 86.1 & 0.118 & 94.3 & 0.049 \\
\midrule
{Ours-1B} 
& \underline{94.6} & \underline{0.058} 
& \underline{99.6} & \underline{0.010} 
& \underline{99.8} & \underline{0.012} 
& \underline{98.4} & \underline{0.038} 
& \underline{89.5} & \underline{0.097} 
& \underline{97.4} & \underline{0.041} \\
{Ours-10B} 
& \textbf{96.3} & \textbf{0.050} 
& \textbf{99.7} & \textbf{0.007} 
& \textbf{99.8} & \textbf{0.009} 
& \textbf{98.7} & \textbf{0.030} 
& \textbf{93.5} & \textbf{0.081} 
& \textbf{98.3} & \textbf{0.035} \\
\bottomrule
\end{tabular}%
}%
\vspace{-10pt}
\end{table*}









An important aspect of \method is to scale the training data, which we achieve by combining a large number of publicly available datasets (\cref{sec:data-sources}) with a new annotation pipeline that we develop to handle dynamic content in off-the-shelf videos (\cref{sec:annotation-pipeline}).

\subsubsection{Data Sources}%
\label{sec:data-sources}

We first collect several publicly available datasets:
Aria series~\cite{pan23aria},
Bedlam~\cite{black23bedlam:},
BEHAVIOR-1K~\cite{li24behavior-1k:},
Co3Dv2~\cite{reizenstein21common},
uCo3D~\cite{liu25uco3d},
DL3DV~\cite{ling23dl3dv-10k:},
Dynamic Replica~\cite{karaev23dynamicstereo},
EDEN~\cite{le21eden:},
EFM3D~\cite{straub24efm3d:},
HOT3D~\cite{banerjee24introducing},
Habitat~\cite{habitat19iccv},
Hypersim~\cite{roberts21hypersim:},
Mapfree~\cite{arnold22map-free},
Mapillary Metropolis~\cite{mapillary25metropolis},
MPSD~\cite{antequera20mapillary},
Megadepth~\cite{li18megadepth:},
Megasynth~\cite{jiang25megasynth:},
Mid-Air~\cite{fonder19mid-air:},
Mvssynth~\cite{huang18deepmvs:},
ParallelDomain-4D~\cite{hoorick24generative},
Replica~\cite{straub19the-replica},
SAIL-VOS~\cite{hu21sail-vos},
ScanNet Series~\cite{dai17scannet:, yeshwanth23scannet:},
TartanAirV2~\cite{wang20tartanair:},
TartanGround~\cite{patel25tartanground:},
Taskonomy~\cite{zamir18taskonomy:},
UnrealStereo4K~\cite{tosi21smd-nets:},
Virtual KITTI~\cite{cabon20virtual},
Waymo~\cite{sun20scalability}, and
WildRGBD~\cite{xia24rgbd}.
We exclude Kubric~\cite{greff22kubric:} and PointOdyssey~\cite{zheng23pointodyssey:} used by VGGT because their background geometry is fake and yields invalid depth.
Additionally, we use several internal datasets, which include artist-created object assets, rigid and dynamic synthetic environments, real-world device captures, and related sources.
For non-synthetic datasets (\eg, those annotated by SfM pipelines), we remove noisy depth values via a multi-view consistency check and discard sequences with too few valid depth pixels.
In total, these datasets contain approximately $3$M sequences, each containing between $10$ and $20{,}000$ images.

\subsubsection{Data Annotation Pipeline}%
\label{sec:annotation-pipeline}

To further expand our training data, we built a large internal video collection of roughly $40$M Internet-style videos.
We first assess each video for suitability for reconstruction, \eg filtering out clips with large watermarks or abrupt shot changes.
Videos that pass this check are used for self-supervised training as described in \cref{sec:self-supervised-training}.
We then introduce a new pipeline to annotate videos with camera parameters and depth maps, for both static and dynamic scenes.
We prioritize annotation quality over quantity and aggressively reject low-quality data.
We also discard depth annotations in regions that are likely to be dynamic.
This may exclude extreme camera motions or highly dynamic scenes, but these are still well represented in the synthetic datasets.
Overall, we obtain a collection of about $200$K dynamic scenes and $600$K static scenes with high-quality camera and depth annotations.

\paragraph{VLM pre-filtering.}

We prompt a Vision-Language Model (VLM) to discard videos that are unlikely to be reconstructed using multi-view geometry.
The VLM classifies $50\%$ of the clips as too difficult to reconstruct, \eg, because they contain multiple clips, extreme motion blur, or heavy overlays or watermarks.
It also classifies $40\%$ of them as reconstructible, but potentially with low accuracy due to insufficient parallax, lack of non-repetitive texture, etc.
The remaining $10\%$ of videos go to the next stage.
In the same pass, the VLM extracts metadata, such as whether the scene appears dynamic, for use by later stages.

\paragraph{Dynamic mask extraction.}

We use Grounding DINO~\cite{liu24grounding} to detect bounding boxes of potentially movable object categories, such as people and cars.
These regions are then excluded from matching, tracking, and verification.

\paragraph{Feature matching and tracking.}

We extract matches and tracks across frames with an ensemble of methods, using
SIFT~\cite{lowe99object},
SuperPoint \& SuperGlue~\cite{daniel18superpoint:, sarlin20superglue:},
ALIKED \& LightGlue~\cite{zhao23aliked:, lindenberger23lightglue:}, and
VGGSfM Tracker~\cite{wang24vggsfm:}.
Matches within dynamic regions are discarded.

\paragraph{Reconstruction and filtering.}

We use the original VGGT to initialize camera parameters when RANSAC-based essential matrix estimation yields too few inliers, and then run COLMAP~\cite{schonberger16structure-from-motion} for iterative bundle adjustment and filtering based on the correspondences computed above.
For successful reconstructions, we discard sequences that fail heuristic checks, \eg, an image registration ratio $<99.5\%$, a field of view outside $[30^\circ,120^\circ]$, or a distortion ratio $>0.1$.
These criteria aggressively remove cases with degenerate motion or extreme zoom.
Then, we estimate per-frame dense depth maps using patch-based multi-view stereo~\cite{schonberger16pixelwise}.

\paragraph{Multi-view consistency.}

For each frame, we unproject the depth map to 3D, reproject the points into other views, and compare them with the depths there.
Pixels that satisfy this cross-view consistency check are marked valid.
We discard sequences with fewer than $5\%$ of pixels with valid depth, which typically, though not always, indicates low-quality cameras.

\paragraph{Supervised geometric filtering.}

Finally, we obtain cameras, depths, and valid masks for every sequence.
We use handcrafted features, such as camera-up-vector consistency, parallax angle, and trajectory smoothness, to describe each sequence.
We hand-annotated 500 static and 500 dynamic sequences, respectively, to train a classifier to remove low-quality reconstructions.
The classifier is an ensemble of XGBoost~\cite{chen16xgboost:}, random forests~\cite{breiman01random}, and CatBoost~\cite{prokhorenkova18catboost:}.
